\item \label{itm:thermtec} {\bf [5 pts.]} Thermotectonic evolution of thrusts.\\
\begin{figure}[hbt]
    \centering
    \includegraphics[]{figures/tectonic_processes.eps}
    \caption{Thermal evolution of several idealized tectonic terranes.}
    \label{fig:thermtec}
\end{figure}

Both continental orogens and accretionary wedges are built structurally by
thrusting, yet the thermal evolution appears quite different.  Consider the
hypothetical thermal evolution of two tectonic environments presented in
Figure~\ref{fig:thermtec}.  A stable craton is provided as a reference. 

Describe the features and change in properties/sources that give rise to the
specific thermal evolution pattern related a thrust (Figure~\ref{fig:thermtec}).
\begin{enumerate}
    \item Why does heat flow initially decrease?
    \item Why does heat flow rebound?
    \item Why does heat flow reach a peak in an orogen, but not an accretionary
    wedge?
    \item Why does the long-term thermal evolution of a thrust take so long to
    approach that of a craton?
    \item Why should cratonic heat flow decrease slightly with time?
\end{enumerate}

As part of your discussion:
\begin{itemize}
    \item Draw the idealized geotherm immediately following the thermotectonic
    event.  Label this curve $t_0$.

    \item Using a different colored pen/pencil, draw the equilibrium geotherm.
    Label this curve $t_\infty$.

    \item Using a third color, draw three hypothetical cooling at a few
    instances in time.  Label these curves $t_1$, $t_2$, and $t_3$.

    \item Your result will invariably look different than the model in
    Figure~\ref{fig:thermtec}.  What additional processes or changes to the
    lithosphere occur during a thrusting that may contribute to this difference?
    Why?
\end{itemize}
\clearpage
